% Elaborado por Fabian Gomez
% Version 1.0

\section{System Requirements}\label{sec:requirements}

Esta sección contiene la ``definición maestra'' de los requisitos del sistema (RF/RNF/CON) del Rover Olympus y sus derivados a nivel de subsistema, organizada según la arquitectura de 7 subsistemas. Los requisitos se consolidan aquí para asegurar la coherencia del baseline técnico. \cite{omg_sysml_mbse, nasa_cm, segoldmine_ppi}

\subsection{Functional Requirements (By Subsystem)}

Los requisitos funcionales describen capacidades obligatorias del Rover Olympus para ejecutar UC‑01. Se organizan bajo los 7 subsistemas técnicos definidos en la arquitectura.

\subsubsection{Structural (Estructura / Chasis)}
\textbf{Soporte de Sistema:}
\begin{itemize}
    \item \textbf{STR-REQ-001:} El subsistema Structural DEBERÁ soportar la masa total del rover manteniendo un factor de seguridad de al menos 1.5 frente a cargas estáticas en pendientes de hasta 20°. [EXTERNO]
    \item \textbf{STR-REQ-002:} El subsistema Structural DEBERÁ proveer una envolvente de protección para la electrónica crítica equivalente a IP4X para prevenir el ingreso de suelo volcánico. [EXTERNO]
\end{itemize}

\subsubsection{Tracción (Locomotion)}
\textbf{Soporte de RF-003, RF-004:}
\begin{itemize}
    \item \textbf{TRAC-REQ-001:} El subsistema de Tracción DEBERÁ permitir el desplazamiento estable en superficies volcánicas con pendientes de hasta 20°.
    \item \textbf{TRAC-REQ-002:} El subsistema de Tracción DEBERÁ proveer un despeje del suelo (ground clearance) de al menos 10 cm. [EXTERNO]
\end{itemize}

\subsubsection{Propulsión (Drive / Actuación)}
\textbf{Soporte de RF-005:}
\begin{itemize}
    \item \textbf{PROP-REQ-001:} El subsistema de Propulsión DEBERÁ generar torque suficiente para iniciar el movimiento desde el reposo en una pendiente de 20°.
    \item \textbf{PROP-REQ-002:} El subsistema de Propulsión DEBERÁ detener el giro de las ruedas en un tiempo $\le$ 500 ms tras una señal de Fault Stop (soporte a RF-005).
\end{itemize}

\subsubsection{C\&DH (Command \& Data Handling)}
\textbf{Soporte de RNF-001, RF-006:}
\begin{itemize}
    % CDH-REQ-001 (abstraído de Yocto)
    \item \textbf{CDH-REQ-001:} El subsistema C\&DH \textbf{DEBERÁ} proveer un entorno de ejecución basado en un sistema operativo de tiempo real o Linux embebido a medida, capaz de sostener el procesamiento de agentes en \(\le 2\,\mathrm{s}\) (soporte a RNF-001).
    \item \textbf{CDH-REQ-002:} El subsistema C\&DH DEBERÁ disponer de almacenamiento local suficiente para 2 horas de telemetría y logs de misión. [EXTERNO]
\end{itemize}

\subsubsection{Energía (EPS / Power)}
\textbf{Soporte de CON-001:}
\begin{itemize}
    \item \textbf{EPS-REQ-001:} El subsistema de Energía DEBERÁ reportar voltaje y corriente al C\&DH con una frecuencia de $\ge$ 1 Hz y precisión de $\pm$ 1\% mediante sensores de monitoreo de potencia integrados.
    \item \textbf{EPS-REQ-002:} El subsistema de Energía DEBERÁ interrumpir la potencia a los motores ante voltajes críticos de la batería principal del sistema. [EXTERNO]
\end{itemize}

\subsubsection{COMM (Telecomunicaciones)}
\textbf{Soporte de RF-006, CON-002:}
\begin{itemize}
    \item \textbf{COMM-REQ-001:} El subsistema COMM DEBERÁ implementar el protocolo CSP para asegurar la integridad de TM/CMD (soporte a RF-006).
    \item \textbf{COMM-REQ-002:} El subsistema COMM DEBERÁ mantener un enlace estable a 2.4 GHz hasta 50 m de distancia en línea de vista. [EXTERNO]
\end{itemize}

% ... (UHF requirements remain focused on protocols as they are interface standards)

\subsubsection{GNC (Guía, Navegación y Control)}
\textbf{Soporte de RF-001, RF-002, RF-003, RF-004:}
\begin{itemize}
    \item \textbf{GNC-REQ-001:} El subsistema GNC DEBERÁ adquirir imágenes mediante una interfaz de cámara de alta velocidad con un FOV horizontal $\ge$ 120° (soporte a RF-001).
    \item \textbf{GNC-REQ-002:} El subsistema GNC DEBERÁ clasificar la transitabilidad con una confianza $\ge$ 95\% en sectores de 10x10 cm (soporte a RF-002).
    \item \textbf{GNC-REQ-003:} El subsistema GNC DEBERÁ ejecutar fusión sensorial IMU+Encoders a una frecuencia mínima de 10 Hz para estimar el slip (soporte a RF-004). [EXTERNO]
\end{itemize}

\subsection{System Functional Requirements (Definición Maestra)}
\textbf{RF-001: Percepción Visual (CSI)} \\
El sistema DEBERÁ adquirir percepción visual mediante cámara CSI frontal. Ver GNC-REQ-001.

\textbf{RF-002: Clasificación de Transitabilidad} \\
El sistema DEBERÁ clasificar la transitabilidad del terreno para navegación autónoma. Ver GNC-REQ-002.

\textbf{RF-003: Evitación Reactiva de Obstáculos} \\
El sistema DEBERÁ evitar obstáculos $\ge$ 20 cm. Ver GNC-REQ-002 y TRAC-REQ-002.

\textbf{RF-004: Compensación de Slip} \\
El sistema DEBERÁ compensar dinámicamente el deslizamiento. Ver GNC-REQ-003.

\textbf{RF-005: Fault Stop Resiliente} \\
El sistema DEBERÁ ejecutar parada segura ante slip $>$ 20\%. Ver PROP-REQ-002.

\textbf{RF-006: Integridad de Telemetría (CSP)} \\
El sistema DEBERÁ asegurar integridad vía CSP. Ver COMM-REQ-001.

\subsection{Performance and Non-Functional Requirements}
\sysreq{RNF-001}{Latencia de Procesamiento}{CRÍTICO}{ELANav}{$\le$ 2 s por ciclo.}
\sysreq{RNF-002}{Resiliencia de Recuperación}{ALTO}{ConOps}{Éxito $\ge$ 90\% en rescates.}
\sysreq{RNF-003}{Precisión de Navegación}{ALTO}{ConOps}{Error $<$ 5\% dist. total.}

\subsection{Design and Implementation Constraints}
\sysreq{CON-001}{Simp. Energética}{BAJA}{ConOps}{Operación Lab-only.}
\sysreq{CON-002}{Inyección de Latencia}{ALTA}{ConOps}{Delay 0-5s simulado.}
